\documentclass[10pt,a4paper]{amsart}
\usepackage[margin=19mm]{geometry}
\usepackage{amsmath,amssymb,mathtools}
\usepackage[scr=rsfs,cal=euler]{mathalfa}
\usepackage{xfrac}
\usepackage[unicode,hidelinks,bookmarks=false]{hyperref}
\newcommand{\ie}{i.e.\ }
\newcommand{\cf}{cf.\ }
\newcommand{\resp}{respectively\ }
\newcommand{\Subsection}{\subsection}
\providecommand{\nobreakdash}{\nobreak\hbox{-}}
\setlength{\emergencystretch}{2em}
\multlinegap0pt
\usepackage{amssymb}
\usepackage{mathtools}
\newtheorem{theorem}{Theorem}
\newtheorem{lemma}{Lemma}
\newcommand{\E}{\mathcal E_p}
\DeclareMathOperator{\Gal}{Gal}
\newcommand{\Hp}{H_{p^3}}
\title{From Common-Slot Chains to Heisenberg Central Products over Global Fields}
\author{Marina Palaisti}
\date{Source: arXiv:2609.00244v1 (CC BY 4.0)}
\begin{document}
\maketitle

\noindent\emph{Source and licence.} From Common-Slot Chains to Heisenberg Central Products over Global Fields. Version arXiv:2609.00244v1 (CC BY 4.0), distributed by arXiv under the Creative Commons Attribution 4.0 International licence, http://creativecommons.org/licenses/by/4.0/, as recorded in the arXiv licence metadata for this exact version and shown on the abstract page https://arxiv.org/abs/2609.00244v1. Full licence notice: \texttt{licenses/arxiv-2609.00244-CC-BY-4.0.txt}.\par\smallskip

\noindent\textit{Editorial scope.} Counted unit: the article's theorem thm:explicit (source0.tex lines 640--673). The article's complete original proof of it is delivered with it (source0.tex lines 675--840, one \texttt{proof} environment of 166 lines). No mathematics is written, repaired or re-derived: the statement and the proof are the article's own lines, in the article's own notation, and no retained line is substituted or re-flowed.  Retained uncounted as the source-local context the proof needs: lines 469--471; lines 537--553; lines 606--611.  Nothing was omitted from any retained span, so the package carries no omission record.  No retained line contains a citation, so the wrapper carries no bibliography and no bibliography entry.  No retained line contains a figure environment, an image-inclusion command or drawing code, so no image is delivered and the package declares no asset dependency.  Editorial formatting only, in addition: an AMS \texttt{amsart} preamble that restates the article's own declarations and macros used by the retained lines, namely the environments \texttt{theorem}, \texttt{lemma} on separate counters, together with the standard packages those lines need.  The retained environments print in this document as Theorem 1, Lemma 1 in order of appearance, whereas the article prints the same items with the numbering of its own full document.  The complete native source of the article as distributed in the arXiv e-print for this exact version is bundled unchanged as \texttt{sources/209007/source0.tex}.
\begin{equation}\label{eq:E-extension}
 1\longrightarrow C_p\longrightarrow\E\longrightarrow C_p^4\longrightarrow1.
\end{equation}

\begin{theorem}\label{thm:global-realizability}
Let $F$ be a global field, let $p$ be an odd prime with
$p\neq\operatorname{char}F$, and assume $\mu_p\subset F$. Let
\[
 K=F(\sqrt[p]{a},\sqrt[p]{b},\sqrt[p]{c},\sqrt[p]{d})
\]
be a $C_p^4$-extension. There exists a Galois extension $L/F$ containing
$K$ such that
\[
 \Gal(L/F)\cong\Hp*\Hp
\]
and the natural map $\Gal(L/F)\to\Gal(K/F)$ is the quotient map in
\eqref{eq:E-extension} if and only if
\[
 (a,b)_p=(c,d)_p.
\]
\end{theorem}

\begin{lemma}\label{lem:W-ratio}
For $f\in F(\sqrt[p]{u})^\times$,
\[
 \frac{\rho_u(W_u(f))}{W_u(f)}=\frac{N_u(f)}{f^p}.
\]
\end{lemma}

\begin{theorem}\label{thm:explicit}
Assume, in addition to the hypotheses of
Theorem~\ref{thm:global-realizability}, that
\[
 (a,b)_p=(c,d)_p
\]
and that the six classes
\[
 a,b,c,d,x,y\in F^\times/F^{\times p}
\]
are linearly independent. Put
\[
 M=F(\sqrt[p]{a},\sqrt[p]{b},\sqrt[p]{c},\sqrt[p]{d},
   \sqrt[p]{x},\sqrt[p]{y}).
\]
Choose $r\in F^\times$ such that $rw\notin M^{\times p}$, put $w'=rw$,
and set
\[
 \widetilde L=M(\sqrt[p]{w'}).
\]
Then
\[
 \Gal(\widetilde L/F)\cong(\Hp*\Hp)\times C_p^2.
\]
If $A\cong C_p^2$ denotes the factor generated by the two auxiliary Kummer
directions $x$ and $y$, then
\[
 L=\widetilde L^A
\]
contains $K$ and satisfies
\[
 \Gal(L/F)\cong\Hp*\Hp.
\]
\end{theorem}

\begin{proof}
First we justify the choice of $r$. By elementary Kummer theory, the kernel
of the natural map
\[
 F^\times/F^{\times p}\longrightarrow M^\times/M^{\times p}
\]
is precisely the six-dimensional subspace generated by the classes of
$a,b,c,d,x,y$. Indeed, if $u\in F^\times$ becomes a $p$-th power in $M$,
then $F(\sqrt[p]{u})$ is a degree-$p$ subextension of $M/F$, and hence its
Kummer class lies in the span of the six defining classes; the converse is
immediate.

Moreover, $F^\times/F^{\times p}$ is infinite for a global field. For
example, given any finite set of classes, choose a finite place outside the
supports of representatives of those classes and use weak approximation to
produce an element having valuation $1$ at that place. Its class cannot lie
in the prescribed finite-dimensional subspace. The set of classes
$[r]\in F^\times/F^{\times p}$ for which $rw\in M^{\times p}$ is either
empty or a single coset of the above kernel: if both $r_1w$ and $r_2w$ are
$p$-th powers in $M$, then $r_1/r_2$ is a $p$-th power in $M$. Hence we may
choose $r\in F^\times$ with $rw\notin M^{\times p}$. Multiplying $w$ by
$r\in F^\times$ does not change any quotient $\sigma(w)/w$.

Let
\[
 z=\sqrt[p]{w'}.
\]
Since $w'\notin M^{\times p}$, the extension $\widetilde L/M$ has degree
$p$. Because $r\in F^\times$, for every $\sigma\in\Gal(M/F)$ we have
\[
 \frac{\sigma(w')}{w'}=\frac{\sigma(w)}{w}.
\]
Because the six Kummer classes are independent,
\[
 \Gal(M/F)=
 \langle\sigma_a,\sigma_b,\sigma_c,\sigma_d,\sigma_x,\sigma_y\rangle
 \cong C_p^6,
\]
where each $\sigma_u$ multiplies $\sqrt[p]{u}$ by $\zeta$ and fixes the
other five chosen radicals.

On $F(\sqrt[p]{by})$, the restrictions of $\sigma_b$ and $\sigma_y$ are both
$\rho_{by}$; similarly, on $F(\sqrt[p]{xc})$, the restrictions of
$\sigma_x$ and $\sigma_c$ are both $\rho_{xc}$. Lemma~\ref{lem:W-ratio}
and the four norm equations therefore give
\begin{align*}
 \frac{\sigma_b(w)}{w}
 &=\frac{ax}{f_1^p}\frac{x^{-1}}{f_2^p}
 =\left(\frac{\sqrt[p]{a}}{f_1f_2}\right)^p,\\
 \frac{\sigma_y(w)}{w}
 &=\frac{x^{-1}}{f_2^p}
 =\left(\frac{1}{\sqrt[p]{x}f_2}\right)^p,\\
 \frac{\sigma_c(w)}{w}
 &=\frac{y^{-1}}{f_3^p}\frac{yd}{f_4^p}
 =\left(\frac{\sqrt[p]{d}}{f_3f_4}\right)^p,\\
 \frac{\sigma_x(w)}{w}
 &=\frac{y^{-1}}{f_3^p}
 =\left(\frac{1}{\sqrt[p]{y}f_3}\right)^p,
\end{align*}
and $\sigma_a(w)=\sigma_d(w)=w$.

Hence the Kummer automorphisms extend to $\widetilde L$ by
\begin{align*}
 \widetilde\sigma_a(z)&=z,
 &\widetilde\sigma_d(z)&=z,\\
 \widetilde\sigma_b(z)&=\frac{\sqrt[p]{a}}{f_1f_2}z,
 &\widetilde\sigma_y(z)&=\frac{1}{\sqrt[p]{x}f_2}z,\\
 \widetilde\sigma_c(z)&=\frac{\sqrt[p]{d}}{f_3f_4}z,
 &\widetilde\sigma_x(z)&=\frac{1}{\sqrt[p]{y}f_3}z.
\end{align*}
Let $\tau$ be the automorphism fixing $M$ and satisfying
\[
 \tau(z)=\zeta z.
\]
These lifts generate $\Gal(\widetilde L/F)$ because their restrictions
generate $\Gal(M/F)$ and they contain the generator $\tau$ of
$\Gal(\widetilde L/M)$.

Each displayed lift has order $p$. For example,
\[
 \prod_{i=0}^{p-1}\sigma_b^i
 \left(\frac{\sqrt[p]{a}}{f_1f_2}\right)
 =\frac{a}{N_b(f_1)N_{by}(f_2)}
 =\frac{a}{(ax)(x^{-1})}=1,
\]
and
\[
 \prod_{i=0}^{p-1}\sigma_y^i
 \left(\frac{1}{\sqrt[p]{x}f_2}\right)
 =\frac{1}{xN_{by}(f_2)}=1.
\]
The $c$- and $x$-cases are identical; the $a$- and $d$-lifts plainly have
order $p$.

For commuting quotient automorphisms $\sigma_i,\sigma_j$ with lifts
\[
 \widetilde\sigma_i(z)=q_i z,
 \qquad
 \widetilde\sigma_j(z)=q_j z,
\]
our commutator convention gives
\begin{equation}\label{eq:comm-formula}
 [\widetilde\sigma_i,\widetilde\sigma_j](z)
 =\frac{\sigma_i(q_j)q_i}{\sigma_j(q_i)q_j}z.
\end{equation}
Using \eqref{eq:comm-formula},
\[
 [\widetilde\sigma_a,\widetilde\sigma_b]=\tau,
\]
because $\sigma_a(\sqrt[p]{a})=\zeta\sqrt[p]{a}$ and $\sigma_a$ fixes
$f_1,f_2$. Likewise,
\[
 [\widetilde\sigma_c,\widetilde\sigma_d]=\tau^{-1},
\]
because $\sigma_d(\sqrt[p]{d})=\zeta\sqrt[p]{d}$.
All other commutators among
\[
 \widetilde\sigma_a,\widetilde\sigma_b,
 \widetilde\sigma_c,\widetilde\sigma_d
\]
are trivial.

It remains to check that the two auxiliary lifts generate a commuting
$C_p^2$-factor relative to these four lifts. All commutators are immediate
from the displayed quotients except for
\[
 (\widetilde\sigma_b,\widetilde\sigma_y),\qquad
 (\widetilde\sigma_c,\widetilde\sigma_x),\qquad
 (\widetilde\sigma_x,\widetilde\sigma_y).
\]
For the first pair, $\sigma_b$ and $\sigma_y$ have the same restriction to
$F(\sqrt[p]{by})$, so the two occurrences of the conjugate of $f_2$ in
\eqref{eq:comm-formula} cancel. The second pair is identical. For the last
pair, writing
\[
 q_y=\frac{1}{\sqrt[p]{x}f_2},
 \qquad
 q_x=\frac{1}{\sqrt[p]{y}f_3},
\]
we have
\[
 \sigma_x(q_y)=\zeta^{-1}q_y,
 \qquad
 \sigma_y(q_x)=\zeta^{-1}q_x,
\]
so the factors again cancel.

Thus
\[
 \langle\widetilde\sigma_a,\widetilde\sigma_b,
 \widetilde\sigma_c,\widetilde\sigma_d,\tau\rangle
 \cong\Hp*\Hp,
\]
while $\widetilde\sigma_x$ and $\widetilde\sigma_y$ generate a commuting
$C_p^2$-factor. Since
\[
 [\widetilde L:F]=p\,[M:F]=p^7=|\Hp*\Hp|\,p^2,
\]
these relations account for the full Galois group, and therefore
\[
 \Gal(\widetilde L/F)\cong(\Hp*\Hp)\times C_p^2.
\]
The fixed field of the auxiliary factor contains $K$, because the auxiliary
generators fix $\sqrt[p]{a},\sqrt[p]{b},\sqrt[p]{c},\sqrt[p]{d}$, and its
Galois group over $F$ is $\Hp*\Hp$.
\end{proof}

\end{document}
